\documentclass[11pt,a4paper]{article}
\usepackage[utf8]{inputenc}
\usepackage[T1]{fontenc}
\usepackage{geometry}
\usepackage{hyperref}
\usepackage{booktabs}
\usepackage{enumitem}
\usepackage{listings}
\usepackage{xcolor}
\geometry{margin=2.5cm}

\lstset{
  basicstyle=\ttfamily\small,
  breaklines=true,
  frame=single,
  backgroundcolor=\color{gray!8},
  numbers=left,
  numberstyle=\tiny
}

\title{02244 Logic for Security\\Project on Security Protocols\\
\vspace{0.5em}\large Open Authorization Protocol}
\author{
\textbf{Authors}\\
Zeinab\\
Fahmida\\
Mirtha\\[0.5em]
March 2026
}
\date{}

\begin{document}
\maketitle
\thispagestyle{empty}

\newpage
\tableofcontents
\newpage

\section{Introduction}

This project designs and analyzes an Open Authorization protocol for the following scenario:
user $A$ stores photos at server $B$ and wants photo-book service $P$ to print selected photos.
Instead of downloading and re-uploading data, $A$ should authorize $P$ to access only the intended subset at $B$.

The trust model includes an identity provider $I$ (IDP), similar to MitID-like systems.
User $A$ authenticates through a shared password relation with $I$, while authorization tokens are issued by $I$ and presented to $B$ through $P$.
The development follows the weekly progression in the course: Dolev--Yao analysis, lazy intruder reasoning,
typing/format discipline, channel abstraction, and guessable-password verification.

This report documents not only the final protocol idea but also intermediate failed or vulnerable versions,
as required by the assignment. Those intermediate steps are essential because they explain
why each refinement was introduced.

\section{Method and Tooling}

We model protocol versions in AnB and verify them with OFMC 2024.
Each week contributes one or more development artifacts under \texttt{ProjectProtocol/}.
For each version, we record:
\begin{itemize}[leftmargin=1.2em]
  \item protocol intent and assumptions,
  \item formal goals,
  \item OFMC output (\texttt{ATTACK\_FOUND} or \texttt{NO\_ATTACK\_FOUND}),
  \item interpretation and next refinement decision.
\end{itemize}

The analysis is symbolic under Dolev--Yao assumptions.
Where bounded-session settings are used, this is explicitly stated.

\section{Protocol Development by Week}

\subsection{Week 2: Bare-Bones Baseline and Static Dolev--Yao}

\textbf{Files used:}
\begin{itemize}[leftmargin=1.2em]
  \item \texttt{Week2\_Problem1\_Informal\_Protocol\_Outline.md}
  \item \texttt{Week2\_Problem2\_First\_AnB\_Version.anb}
  \item \texttt{Week2\_Problem3\_Static\_DolevYao\_Analysis.md}
\end{itemize}

The Week 2 protocol is intentionally minimal:
$A$ signs a delegation context and forwards authorization for $P$ to access $B$.
This version is used to establish a runnable baseline and perform the required static passive-intruder analysis.

\textbf{Main finding:}
the static Dolev--Yao derivation shows that confidentiality can fail in the baseline model under intruder-controlled delegation behavior.
This motivates introducing explicit IDP issuance in Week 3.

\subsection{Week 3: IDP-Based Delegation and Lazy Intruder Task}

\textbf{Files used:}
\begin{itemize}[leftmargin=1.2em]
  \item \texttt{Week3\_Problem1\_AnB\_With\_IdentityProvider.anb}
  \item \texttt{Week3\_Problem2\_LazyIntruder\_Executability\_of\_P.md}
  \item \texttt{Week3\_Problem3\_Refined\_HonestProvider.anb}
\end{itemize}

Week 3 introduces role $I$ as the token issuer.
User $A$ authenticates to $I$ with a password-derived hash witness,
and $I$ signs a grant that is forwarded via $P$ to $B$.
This satisfies the requirement that password is not used as an encryption key.

\textbf{Main findings:}
\begin{itemize}[leftmargin=1.2em]
  \item The special task is completed: role $P$ executability is shown step-by-step with lazy intruder method.
  \item OFMC still reports attacks in this stage (\texttt{ATTACK\_FOUND}), showing refinement is still needed.
\end{itemize}

\subsection{Week 4: Formats, Type-Flaw Resistance, and Key Lookup}

\textbf{Files used:}
\begin{itemize}[leftmargin=1.2em]
  \item \texttt{Week4\_Problem1\_MainProtocol\_Formats.anb}
  \item \texttt{Week4\_Problem2\_AnB\_KeyLookup\_Protocol.anb}
  \item \texttt{Week4\_Problem3\_TypeFlawResistance\_Check.md}
\end{itemize}

Week 4 focuses on message structure and implementation realism.
We introduced explicit format constructors
(\texttt{fReq}, \texttt{fAuthProof}, \texttt{fGrant}, \texttt{fUse}, \texttt{fData})
to reduce ambiguous unification patterns.
We also separated public-key retrieval into a dedicated protocol where $A$ asks $I$ for $pk(X)$.

\textbf{Main findings:}
\begin{itemize}[leftmargin=1.2em]
  \item Several intermediate formatted variants were non-executable in OFMC and had to be simplified.
  \item Final Week 4 models are executable and suitable for type-flaw discussion.
  \item Weak-auth attacks remain in this stage.
\end{itemize}

\subsection{Week 5: Pseudonymous Channels (TLS-Style Abstraction)}

\textbf{Files used:}
\begin{itemize}[leftmargin=1.2em]
  \item \texttt{Week5\_Problem1\_PseudonymousChannels.anb}
  \item \texttt{Week5\_Problem1\_PseudonymousChannels\_Result.md}
\end{itemize}

This week replaces explicit transport cryptography with pseudonymous channel notation,
as required by the channel-abstraction task.
The goal is to keep authorization logic while delegating transport assumptions to channel semantics.

\textbf{Main finding:}
the model runs, but OFMC reports \texttt{ATTACK\_FOUND} on secrecy (\texttt{fData(Photos)} leak).
Hence channel abstraction alone does not guarantee protocol security.

\subsection{Week 6: Guessable Password Verification}

\textbf{Files used:}
\begin{itemize}[leftmargin=1.2em]
  \item \texttt{Week6\_Problem1\_GuessablePassword\_Check.anb}
  \item \texttt{Week6\_Problem2\_GuessablePassword\_Hardened.anb}
  \item \texttt{Week6\_Problem3\_GuessablePassword\_OnlyCheck.anb}
  \item \texttt{Week6\_Problem1\_Result\_Summary.md}
\end{itemize}

Week 6 addresses low-entropy password security directly:
\begin{enumerate}[leftmargin=1.4em]
  \item Baseline guessable-password model reports \texttt{ATTACK\_FOUND} with \texttt{guesswhat}.
  \item Hardened model (extra fresh value in password proof) shifts first failure away from \texttt{guesswhat}.
  \item Focused guessable-only check reports \texttt{NO\_ATTACK\_FOUND} for bounded run (\texttt{--numSess 1}).
\end{enumerate}

This is the key Week 6 conclusion: for the dedicated special-task goal and tested bound,
guessable-password attack is not found in the focused hardened check.

\section{Special Tasks: Direct Answers}

\begin{tabular}{@{}p{1.5cm}p{5cm}p{7.5cm}@{}}
\toprule
Week & Task & Answer in this project \\
\midrule
2 & Static Dolev--Yao analysis & Completed in \texttt{Week2\_Problem3\_Static\_DolevYao\_Analysis.md}; baseline weakness identified. \\
3 & Lazy intruder executability of role $P$ & Completed in \texttt{Week3\_Problem2\_LazyIntruder\_Executability\_of\_P.md} with lecture-style derivation. \\
4 & Type-flaw resistance + key lookup & Formats and argument in Week 4 files; separate key lookup protocol provided. \\
5 & Pseudonymous channel variant & Implemented in \texttt{Week5\_Problem1\_PseudonymousChannels.anb}; analyzed with OFMC. \\
6 & Guessable password robustness & Progressive check completed; final focused bounded run gives \texttt{NO\_ATTACK\_FOUND}. \\
\bottomrule
\end{tabular}

\section{Final Protocol Status and Limitations}

The project demonstrates a complete formal development process rather than a single final ``perfect'' protocol.
Several variants remain vulnerable to secrecy or weak-auth attacks under OFMC, and these are intentionally retained in the logbook evidence.
This is consistent with the assignment requirement to document both successful and unsuccessful development stages.

Main remaining limitations:
\begin{itemize}[leftmargin=1.2em]
  \item Some variants satisfy specific goals while failing others.
  \item Bounded-session results do not imply full unbounded security.
  \item Channel abstraction and typing discipline improve structure but do not automatically solve authorization logic flaws.
\end{itemize}

\section{Reproducibility Commands}
\begin{lstlisting}
ofmc --numSess 2 "ProjectProtocol\Week2_Problem2_First_AnB_Version.anb"
ofmc --numSess 2 "ProjectProtocol\Week3_Problem1_AnB_With_IdentityProvider.anb"
ofmc --numSess 2 "ProjectProtocol\Week3_Problem3_Refined_HonestProvider.anb"
ofmc --numSess 2 "ProjectProtocol\Week4_Problem1_MainProtocol_Formats.anb"
ofmc --numSess 2 "ProjectProtocol\Week4_Problem2_AnB_KeyLookup_Protocol.anb"
ofmc --numSess 2 "ProjectProtocol\Week5_Problem1_PseudonymousChannels.anb"
ofmc --numSess 2 "ProjectProtocol\Week6_Problem1_GuessablePassword_Check.anb"
ofmc --numSess 2 "ProjectProtocol\Week6_Problem2_GuessablePassword_Hardened.anb"
ofmc --numSess 1 "ProjectProtocol\Week6_Problem3_GuessablePassword_OnlyCheck.anb"
\end{lstlisting}

\section*{Resources and AI Disclosure}
\begin{itemize}[leftmargin=1.2em]
  \item Assignment handout: \texttt{protocol-project (2).pdf}
  \item Course lecture material (Weeks 2--6 topics)
  \item OFMC 2024 and AnB modeling
  \item AI support used for drafting and restructuring report/logbook text;
  all security claims validated against observed OFMC outputs.
\end{itemize}

\end{document}
